% TOP_PROBLEM: {"id":30007645,"problem_number":"LOCAL-30007645","title":"Causal channels of family advantage","category_id":21,"category":{"id":21,"name":"economics","display_name":"Economics","description":"Open economic theory statements, published research agendas, and proposed scoped specifications; question provenance and historical priority are qualified per record.","slug":"economics","order_index":21,"created_at":"2026-10-08T00:00:00Z"},"status":"provenance_pending","external_url":null,"legacy_ids":[],"published":false,"statement":"In the explicitly specified setting: How much intergenerational persistence comes from parental financial resources rather than skills, preferences, networks, assortative mating, and shared environments? The mathematical statement above fixes the target, assumptions and scope of this catalogue formulation.","economics":{"original_id":134,"field":"Inequality and economic mobility","kind":"identification","formalization_basis":"proposed_specification","source_status":"not_verified","priority_status":"not_established","priority_note":"No qualifying problem statement was directly verified, so historical priority cannot be assigned.","earliest_verified_reference":null,"current_reference":{"authors":["Magne Mogstad","Gaute Torsvik"],"title":"Family Background, Neighborhoods and Intergenerational Mobility","year":2021,"url":"https://www.nber.org/papers/w28874","doi":"10.3386/w28874","locator":"Abstract, March 2022 revision, discussion of family environment, genetics, and causal family characteristics","evidence_summary":"The review assesses identification of family channels; the complete proposed causal decomposition is not explicitly certified open in the inspected abstract."},"formalization_note":"This is a proposed scoped research specification. No primary passage explicitly stating this entire remaining question as open was verified. The supporting literature may answer important subquestions; this entry must not be advertised as a verified formal open problem."},"statusQualification":"Provenance pending: an explicit published statement of this exact open question was not verified. This is a catalogue-authored candidate specification, not a certified open conjecture. This is a proposed scoped research specification. No primary passage explicitly stating this entire remaining question as open was verified. The supporting literature may answer important subquestions; this entry must not be advertised as a verified formal open problem.","statusReviewedAt":"2026-10-08"}
% FIRST_POSTED: 2026-10-08
% ENTRY_KIND: statement_only
% !TeX program = lualatex
\documentclass[11pt]{article}
\usepackage{fontspec}
\usepackage[a4paper,margin=25mm]{geometry}
\usepackage{amsmath,amssymb,mathtools}
\usepackage{xurl}
\usepackage[hidelinks]{hyperref}
\newcommand{\sref}[2]{\hyperlink{source:#1:#2}{[S#2]}}
\setlength{\emergencystretch}{3em}
\begin{document}
% problemId:30007645
\section*{Causal channels of family advantage}\label{30007645}
\subsection{Definitions and mathematical statement}
Fix a birth cohort, a parental-income range, and a measure \(Y_i\) of adult economic status. Let \(R_i\) denote parental financial resources during childhood and \(H_i\) parental skills, preferences, and information. Let \(N_i\) represent family networks and \(E_i\) the child's institutional environment. Potential outcomes \(Y_i(r,h,n,e)\) are defined only for declared feasible interventions; they do not imply that every inherited characteristic can be manipulated. Holding prereform \(H_i,N_i,E_i\) fixed defines a controlled resource contrast rather than a total effect that includes their responses. The resource effect is \(\tau_R(r,r';x)=E[Y_i(r,H_i,N_i,E_i)-Y_i(r',H_i,N_i,E_i)\mid X_i=x]\), where \(X_i\) is measured before intervention. Analogous contrasts define network or information interventions. Identify these effects using credible income shocks, randomized supports, adoption or family comparisons with explicit assumptions, and longitudinal measurement. Determine how parental resources interact with skills and constraints rather than assigning all parent--child correlation to money. Mediation claims require additional assumptions beyond total-effect identification. Specify exposure timing, selection into parenthood, and whether parental responses are included. The proposed research task is a context-specific causal decomposition and its policy implications; the cited reviews are not evidence that all mechanisms lack established effects or that one universal decomposition exists.

\par\textbf{Formulation and scope.} This is a proposed scoped research specification. No primary passage explicitly stating this entire remaining question as open was verified. The supporting literature may answer important subquestions; this entry must not be advertised as a verified formal open problem.

\par\textbf{Open-question provenance.} An explicit published open-question statement was not verified. Sources below are supporting literature and must not be treated as a first listing.

\par\textbf{Historical priority.} No qualifying problem statement was directly verified, so historical priority cannot be assigned.

\subsection{Short English statement}
In the explicitly specified setting: How much intergenerational persistence comes from parental financial resources rather than skills, preferences, networks, assortative mating, and shared environments? The mathematical statement above fixes the target, assumptions and scope of this catalogue formulation.

\subsection{Sources}
\begin{itemize}\raggedright
\item[\textnormal{[S1]}]\hypertarget{source:30007645:1}{} Magne Mogstad, Gaute Torsvik. 2021. Family Background, Neighborhoods and Intergenerational Mobility. Abstract, March 2022 revision, discussion of family environment, genetics, and causal family characteristics.
\url{https://www.nber.org/papers/w28874}.
DOI: 10.3386/w28874.
{\small\textit{Supporting or current literature; see evidence qualification. The review assesses identification of family channels; the complete proposed causal decomposition is not explicitly certified open in the inspected abstract.}}
\end{itemize}
\end{document}
